\section{Conclusion}\label{sec:concl}

Lower bounds that scale to large correlation clustering instances used to
certify only factors of about two.  Three observations change this on sparse
graphs: the metric LP and its strengthenings can be restricted to the pairs
at distance at most two; dual block-coordinate ascent turns the restricted
relaxation into an anytime bound whose state is linear in the support; and
star packings on the sparse support are a cheap and strong starting point.
The resulting bounds are exported as certificates that an independent
checker verifies against the raw instance in exact arithmetic.  The gap
decomposition turns the same certificate into a map of where a clustering
can still improve.

\paragraph{Limitations.}
\begin{itemize}
\item On dense instances the root star-packing bound of the KaPoCE
branch-and-bound remains stronger than ours on average; there the block LPs
are too large to converge within the budget.
\item On \numPackBetter{} of the \numSnapCert{} SNAP graphs the sparse star
packing alone gives a larger bound than CertiFlip's combination of a short
packing and block LPs; there the LP part of the method does not pay for its
time within the budget (Section~\ref{sec:exp-cert}).
\item The support has $O(\sum_v\deg(v)^2)$ pairs.  On graphs with strong
hubs it does not fit into memory and no bound is computed.  Restricting the
support to centres of bounded degree would keep the bound valid.
\item The checker is short but not formally verified; the formal development
covers the inequality it evaluates, not its code (Appendix~\ref{app:lean}).
\item The a-priori guarantees of our primal methods are inherited from Pivot,
and the $2.06$ guarantee applies only when one block covers the graph and
separation converges.
\end{itemize}

\paragraph{Open directions.}  Parallel block solves; cluster-LP columns
inside blocks, to which the roundings of \citet{CaoCLLNV24} and
\citet{GarciaSorianoS26} could be applied; overlapping blocks with Lagrangian
cost splitting; and a verified checker, for example by extracting it from the
Lean development.
